\documentclass[11pt]{article}

\usepackage[margin=1in]{geometry}
\usepackage{amsmath,amssymb,amsthm}
\usepackage{booktabs}
\usepackage{enumitem}
\usepackage{hyperref}
\hypersetup{colorlinks=true,linkcolor=blue,citecolor=blue,urlcolor=blue}

\newtheorem{definition}{Definition}
\newtheorem{proposition}{Proposition}
\newtheorem{remark}{Remark}

\DeclareMathOperator*{\argmax}{arg\,max}
\DeclareMathOperator*{\argmin}{arg\,min}
\newcommand{\R}{\mathbb{R}}
\newcommand{\E}{\mathbb{E}}
\newcommand{\PR}{\mathbb{P}}
\newcommand{\Xcal}{\mathcal{X}}
\newcommand{\Hcal}{\mathcal{H}}
\newcommand{\Scal}{\mathcal{S}}
\newcommand{\Gcal}{\mathcal{G}}
\newcommand{\MMD}{\mathrm{MMD}}
\newcommand{\indep}{\perp\!\!\!\perp}

\title{\textbf{Two-Stage Feature Selection and Localization for\\
Distribution-Shift Attribution in a Marketplace}}
\author{Technical write-up}
\date{}

\begin{document}
\maketitle

\begin{abstract}
We formalize a pipeline that, given a new batch of relational marketplace data,
(i) \emph{detects} whether the distribution has shifted, (ii) \emph{localizes}
the shift to the responsible features through a two-stage feature-selection
procedure with hierarchical error control, along two entity axes (supply and
demand), and (iii) turns the result into a structured, machine-consumable
insight carrying a \emph{retrain decision}. Detection uses an exact
permutation Maximum Mean Discrepancy (MMD) test for covariate shift and a
pseudo-outcome / $R$-risk test for concept drift. Localization combines
leave-one-group-out (LOGO) and exact Shapley attribution over the group power
set with a Westfall--Young step-down max-T permutation procedure (family-wise
error control, robust to correlation) and subsampling stability selection. The decision layer estimates covariate and
concept impact on a downstream loss via propensity-based importance weighting,
gated by a batch-separability (AUC) overlap diagnostic. We give definitions,
estimators, validity statements, algorithms, and the structured-insight schema.
\end{abstract}

\tableofcontents

\section{Problem setup and notation}

A marketplace generates \emph{orders} (observations). Each order $o$ records a
buyer (user) $u(o)$, an item $i(o)$, and its merchant $m(o)$ (via
$\text{item}\to\text{merchant}$), an order identifier, and raw numeric
attributes $a(o) \in \R^{A}$ (e.g.\ GMV, quantity, discount, rating). We compare
an \emph{existing} batch (label $W=0$) against a \emph{new} batch ($W=1$).

The shift may be examined along different \emph{entity axes}
$\mathsf{e} \in \{\text{merchant}, \text{buyer}, \dots\}$. Fix an axis with key
$k_{\mathsf e}(o)$ and batch label $W_{\mathsf e}$ defined at the entity level.
Aggregation rolls order-level attributes up to the entity dimension.

\begin{definition}[Aggregation map]
For entity $j$ let $O_j=\{o : k_{\mathsf e}(o)=j\}$ be its orders. The feature map
$\Phi$ produces $x_j=\Phi(O_j)\in\R^{p}$ by applying a fixed set of aggregators
to each raw attribute: summary statistics
$\{\text{mean},\text{sd},\min,\max,\text{median},q_{25},q_{75},\text{sum}\}$,
rolling-window statistics over the entity's order stream ordered by order id, and
activity counts. Feature names factor as
$\texttt{attr}\,\texttt{\_\_}\,\texttt{agg}$, inducing a grouping
$g:\{1,\dots,p\}\to\Gcal$ by raw attribute.
\end{definition}

Let $P_0,P_1$ be the distributions of $x$ on the two batches, with samples
$X=\{x_i\}_{i:W_i=0}$ ($m$ points) and $Y=\{y_i\}_{i:W_i=1}$ ($n$ points). The
global null is $H_0:P_0=P_1$.

\section{Detection}

\subsection{Maximum Mean Discrepancy}

Let $k:\Xcal\times\Xcal\to\R$ be a characteristic kernel (we use the RBF kernel
$k(x,x')=\exp(-\gamma\|x-x'\|^2)$ with $\gamma$ from the median heuristic on the
pooled sample) with RKHS $\Hcal$ and mean embedding $\mu_P=\E_{x\sim P}[k(x,\cdot)]$.

\begin{definition}[MMD]
$\MMD^2(P_0,P_1)=\|\mu_{P_0}-\mu_{P_1}\|_\Hcal^2
=\E[k(x,x')]+\E[k(y,y')]-2\,\E[k(x,y)]$, with $x,x'\sim P_0$, $y,y'\sim P_1$
independent. For characteristic $k$, $\MMD^2(P_0,P_1)=0 \iff P_0=P_1$.
\end{definition}

\begin{definition}[Unbiased estimator]
\begin{equation}
\widehat{\MMD}^2_u
= \frac{1}{m(m-1)}\!\sum_{i\neq j}\! k(x_i,x_j)
+ \frac{1}{n(n-1)}\!\sum_{i\neq j}\! k(y_i,y_j)
- \frac{2}{mn}\!\sum_{i,j}\! k(x_i,y_j).
\label{eq:mmdu}
\end{equation}
\end{definition}
The diagonal is excluded so that $k(x_i,x_i)=1$ does not bias the within-sample
terms; this fact is central to the bootstrap caveat in
Section~\ref{sec:stability}.

\subsection{Permutation test and validity}

Given the pooled sample $Z=X\cup Y$ with labels, draw permutations
$\pi$ of the labels and recompute the statistic. The $p$-value is
$\hat p=(1+\sum_{b=1}^{B}\mathbb{1}\{T_b\ge T_{\mathrm{obs}}\})/(1+B)$.

\begin{proposition}[Exact validity]
Under $H_0$ the pooled points are exchangeable, so for any statistic $T$ the
permutation $p$-value satisfies $\PR_{H_0}(\hat p\le\alpha)\le\alpha$ for all
$\alpha\in(0,1)$, exactly and without distributional assumptions.
\end{proposition}

\subsection{The $k$-sample case}

For $\kappa>2$ batches $P_1,\dots,P_\kappa$, testing $H_0:P_1=\cdots=P_\kappa$
requires aggregating pairwise divergences. A natural omnibus (the kernel analogue
of DISCO / energy-distance ANOVA) is
\begin{equation}
T_{\text{omni}}=\sum_{r<s}\frac{n_r n_s}{N}\,\widehat{\MMD}^2_u(P_r,P_s),
\qquad N=\textstyle\sum_r n_r,
\end{equation}
with a $\kappa$-way label permutation null. A maximum variant
$\max_{r<s}\widehat{\MMD}^2_u$ targets a single divergent pair. The composite
null means rejection must be followed by pairwise/subset localization
(Section~\ref{sec:localization}); this is exactly the two-stage logic.

\subsection{Concept drift via the causal objective}

Detection above concerns $P(X)$. To detect a change in the outcome mechanism
$P(Y\mid X)$ we reuse the project's causal-objective machinery. With batch as
``treatment'' $W$, outcome $Y$, nuisances $\mu(x)=\E[Y\mid x]$ and
$e(x)=\PR(W=1\mid x)$ estimated by cross-fitting, define the pseudo-outcome
$\psi=(Y-\mu(X))(W-e(X))$ and $\tau=\argmin_f \E[(\psi-f(X))^2]$. The PO-risk
statistic $\E[\tau(X)^2]$ (and the $R$-risk analogue) is large under concept
drift and \emph{insensitive to pure covariate shift}; its permute-then-refit
$p$-value gives the concept-drift test. Thus MMD and PO-risk are complementary
detectors, not substitutes.

\section{Localization: attribution by ablation and by coalitions}
\label{sec:localization}

Write $\widehat{\MMD}^2_u(S)$ for the statistic computed on the sub-vector of
features indexed by a subset $S$ of the groups $\Gcal$ (fixed kernel; only the
\emph{ranking} of importances is used, so scale is immaterial).

\begin{definition}[LOGO-MMD]
The leave-one-group-out importance of group $g$ is
$\phi^{\mathrm{LOGO}}_g=\widehat{\MMD}^2_u(\Gcal)-\widehat{\MMD}^2_u(\Gcal\setminus\{g\})$.
\end{definition}
LOGO is the single marginal contribution \emph{at the full coalition}. Its blind
spot is redundancy: two duplicated shifted groups each obtain
$\phi^{\mathrm{LOGO}}\approx 0$.

\begin{definition}[Shapley-MMD]
Treating $v(S)=\widehat{\MMD}^2_u(S)$ as a coalitional game, the Shapley value of
group $g$ is
\begin{equation}
\phi^{\mathrm{Sh}}_g=\sum_{S\subseteq\Gcal\setminus\{g\}}
\frac{|S|!\,(K-|S|-1)!}{K!}\bigl(v(S\cup\{g\})-v(S)\bigr),\quad K=|\Gcal|.
\end{equation}
\end{definition}

\begin{proposition}[Efficiency and exact enumeration]
$\sum_{g}\phi^{\mathrm{Sh}}_g=v(\Gcal)-v(\emptyset)=\widehat{\MMD}^2_u(\Gcal)$,
so the Shapley values additively decompose the total divergence. With
$K$ groups there are $2^{K}$ subsets; for the small $K$ here ($K\!\approx\!7$)
the game is enumerated exactly, requiring no approximation and no gatekeeping.
\end{proposition}
Standalone MMD $v(\{g\})$ (marginal at $\emptyset$), LOGO (marginal at
$\Gcal$), and Shapley (average over all $S$) form a spectrum; unlike LOGO,
$\MMD^2$ is not additive across groups, which is why Shapley---satisfying
efficiency---is the principled decomposition and LOGO the cheap ranking.

\begin{remark}[LOGO $\equiv$ LOCO in a different metric]
The concept-drift analogue replaces the divergence $v$ by a causal risk:
leave-one-covariate-out (LOCO) importance
$\rho_g=\text{risk}(\Gcal\setminus\{g\})-\text{risk}(\Gcal)$ with the PO-/$R$-risk.
The same ablation principle instantiated on $\MMD$ (marginal, covariate) and on
risk (conditional, outcome) covers the two shift types.
\end{remark}

\section{Two-stage selection by permutation and stability}

We deliberately avoid the Benjamini--Hochberg (BH) family. In this problem the
number of group hypotheses is small ($K\!\approx\!7$) and the group statistics
are strongly correlated (aggregations of the same raw attribute, and correlated
raw attributes); BH's independence/PRDS assumptions are dubious and, with so few
hypotheses, BH is simultaneously under-powered and poorly calibrated in
practice. We use permutation-based control instead, which is exact under
exchangeability and correlation-robust.

\paragraph{Stage 1 (groups).} Rank groups by Shapley-MMD; select
$\Gcal^\star=\{g: p_g<\alpha_1 \text{ and } \hat\pi_g\ge\pi_{\mathrm{thr}}\}$,
where $p_g$ is the group's permutation $p$-value and $\hat\pi_g$ its subsampling
selection frequency (Section~\ref{sec:stability}). Combining a permutation test
with a stability threshold is both practical and finite-sample controlled.

\paragraph{Stage 2 (features within a selected group).}
Within each $g\in\Gcal^\star$ apply a \emph{Westfall--Young step-down max-T}
permutation procedure. Let $T_j=\widehat{\MMD}^2_u(\{j\})$ for feature $j\in g$.
Draw $B$ label permutations \emph{shared across the features of $g$}; from each
permutation $b$ form the maximum $M_b=\max_{j\in g}T_j^{(b)}$. The adjusted
$p$-value is
\begin{equation}
\tilde p_j=\frac{1+\sum_{b=1}^{B}\mathbb{1}\{M_b\ge T_j\}}{1+B},
\end{equation}
and feature $j$ is selected if $\tilde p_j\le\alpha_2$.

\begin{proposition}[Exact FWER control]
Under the family null and exchangeability, using the shared-permutation max
statistic gives $\PR(\exists j\ \text{null}: \tilde p_j\le\alpha_2)\le\alpha_2$,
exactly and without independence assumptions; because the null is built from the
same permutations for all features, the procedure adapts to their correlation.
\end{proposition}

The same construction applies at every split of the multi-resolution tree
$\text{root}\to\text{attribute}\to\text{aggregation family}\to\text{feature}$:
each node's children are tested by shared-permutation max-T, and the tree is
\emph{gated}---a node is entered only if its parent was selected---so the
family-wise guarantees compose down a hierarchical, correlation-robust,
BH-free procedure. Each node is re-tested on its own feature subset (marginal at
the root, increasingly conditional deeper down): \emph{subset post-hoc
localization} conditional on a significant global test.

\section{Stability selection and the MMD bootstrap}
\label{sec:stability}

To quantify selection stability we resample and record the selection frequency
$\hat\pi_g$ and a percentile interval of $\phi_g$; a group is stable if
$\hat\pi_g\ge\pi_{\mathrm{thr}}$ and its interval excludes $0$.

\begin{remark}[Bootstrap ties bias the U-statistic]
With-replacement resampling produces distinct positions $i\neq j$ mapping to the
same original point, for which $k(x_i,x_j)=1$ enters the off-diagonal sums in
\eqref{eq:mmdu}. In a resample of size $m$ the expected number of such collisions
is $\approx m-1$, inflating each within-sample term by $\approx(1-c)/(m-1)$ while
the cross term is unaffected; hence $\widehat{\MMD}^2_u$ is biased \emph{upward},
and the effect dominates near $H_0$ where the true value is $\approx0$.
Remedies: additive jitter (a form of kernel smoothing / randomized smoothing),
$m$-out-of-$n$ subsampling without replacement (no ties), generalized diagonal
exclusion of same-origin pairs, or the wild bootstrap for degenerate kernel
statistics. We use the permutation test for \emph{significance} (no ties) and the
bootstrap only for stability of already-significant groups, where these effects
are second order.
\end{remark}

\section{Two-dimensional structure}

Running the pipeline on two axes (merchant and buyer) yields two feature
libraries from the \emph{same} order table. If the batch labels
$W_{\text{merchant}},W_{\text{buyer}}$ are (near-)independent and the injected
shifts act on disjoint raw attributes, aggregating along one axis averages out
the other axis's shift, so each axis recovers only its own drivers. This
separates a \emph{supply-side} change from a \emph{demand-side} change.

\section{Characterization and the retrain decision}

\subsection{Per-feature characterization}
For a selected feature we report the standardized effect
$d=(\bar x_1-\bar x_0)/s_{\mathrm{pool}}$, per-batch quantile grids, and the
KS-optimal split $t^\star=\argmax_t|\hat F_1(t)-\hat F_0(t)|$ with its KS
statistic and per-batch mass below $t^\star$; $t^\star$ materializes into an
executable \texttt{candidate\_rule}.

\subsection{Impact decomposition}
Let a downstream model $\hat f$ be trained on the old batch with per-point loss
$\ell$. Define the in-distribution loss $L_{\mathrm{old}}=\E_{P_0}[\ell]$, the new
loss $L_{\mathrm{new}}=\E_{P_1}[\ell]$, and the importance-weighted estimate
\begin{equation}
L_{\mathrm{iw}}=\E_{P_0}\!\bigl[w(x)\,\ell(x)\bigr],\qquad
w(x)=\frac{p_1(x)}{p_0(x)}=\frac{e(x)}{1-e(x)}\cdot\frac{\PR(W=0)}{\PR(W=1)} ,
\end{equation}
using the propensity $e$ as the density ratio.
\begin{proposition}[Covariate/concept decomposition]
If $P(Y\mid X)$ is unchanged (pure covariate shift) and the supports overlap,
then $L_{\mathrm{iw}}=L_{\mathrm{new}}$, so
\emph{covariate impact} $:=L_{\mathrm{iw}}-L_{\mathrm{old}}$ captures the
degradation and \emph{concept impact} $:=L_{\mathrm{new}}-L_{\mathrm{iw}}=0$.
A significant concept impact therefore indicates concept drift.
\end{proposition}

\subsection{Overlap gate}
The identity requires overlap. Two diagnostics: the weight effective sample size
$\mathrm{ESS}=(\sum w)^2/\sum w^2$, and---robust to weight clipping---the
propensity AUC (batch separability). When $\mathrm{AUC}\approx1$ the batches are
(near-)separable, $w$ is ill-defined, and the covariate/concept split is not
identifiable label-free.

\subsection{Decision rule}
\begin{equation}
\text{decision}=
\begin{cases}
\text{retrain / collect labels}, & \mathrm{AUC}>\tau_{\mathrm{AUC}}\text{ and }
   (L_{\mathrm{new}}-L_{\mathrm{old}})/L_{\mathrm{old}}>\delta,\\
\text{retrain}, & \text{concept-impact CI}>0,\\
\text{reweight}, & (L_{\mathrm{iw}}-L_{\mathrm{old}})/L_{\mathrm{old}}>\delta,\\
\text{monitor}, & \text{otherwise.}
\end{cases}
\end{equation}
Aggregation typically makes a detectable shift separable ($\mathrm{AUC}\to1$), so
at the entity level the honest verdict is usually ``retrain / collect labels,''
with localization indicating \emph{what} to label.

\section{Structured insight contract}

Each finding is emitted as a typed object (\texttt{fsds-insight/1.0}):
identity (\texttt{schema\_version}, \texttt{insight\_id}$=$hash(\texttt{dedup\_key}),
\texttt{dedup\_key}$=$\texttt{axis:scope:target}, \texttt{generated\_at}); a
three-axis confidence $\{\text{significance }p,\ \text{stability }\hat\pi,\
\text{identifiability AUC}\}$; a \texttt{trend}
(\texttt{status}$\in\{$new, worsening, improving, stable, resolved$\}$,
$\Delta$priority, \texttt{first\_seen}) obtained by diffing against a previous
run keyed on \texttt{dedup\_key}; per-feature \texttt{candidate\_rule}; and
routing/ranking (\texttt{owner}, $\text{priority}=\text{exposure}\times\text{impact}$,
\texttt{decision}). The stable identity supports upsert/dedup; the same JSON
drives dashboards, alert routing, and feature-store governance.

\section{Pipeline}

\begin{enumerate}[leftmargin=2em]
\item Aggregate orders to the entity axis: $x_j=\Phi(O_j)$.
\item Global detection: MMD permutation test (covariate); PO-/$R$-risk (concept).
\item Stage 1: Shapley/LOGO ranking; group permutation $p_g$ and stability
      $\hat\pi_g$ $\to$ $\Gcal^\star=\{p_g<\alpha_1,\ \hat\pi_g\ge\pi_{\mathrm{thr}}\}$.
\item Stage 2: within each $g\in\Gcal^\star$, Westfall--Young max-T permutation
      at $\alpha_2$ $\to$ selected features.
\item Characterize each selected feature ($d$, quantiles, KS split, rule).
\item Decision: train downstream model on old batch; compute
      $L_{\mathrm{old}},L_{\mathrm{iw}},L_{\mathrm{new}}$, ESS, AUC $\to$ decision.
\item Emit \texttt{fsds-insight/1.0} objects (two-level: group headline $+$ feature
      drill-down), diff against previous run for \texttt{trend}.
\end{enumerate}

\section{Synthetic validation}

On a synthetic marketplace (250 merchants, 800 buyers, 20{,}000 orders, 69
features/axis) with covariate shift injected on $\{$gmv, user\_rating$\}$
(merchant axis) and $\{$basket\_size, discount\_rate$\}$ (buyer axis):

\begin{center}
\begin{tabular}{lccc}
\toprule
Axis & Global $p$ & Recovered drivers & Leaf false discoveries \\
\midrule
merchant ($W_{\text{merchant}}$) & $0.003$ & $\{$gmv, user\_rating$\}$ & $0$ \\
buyer ($W_{\text{buyer}}$) & $0.003$ & $\{$basket\_size, discount\_rate$\}$ & $0$ \\
\bottomrule
\end{tabular}
\end{center}

Exact Shapley-MMD ranks the true drivers top-2 on each axis (permutation
$p=0.0099$, bootstrap top-2 frequency $1.00$). The decision self-check (with
support overlap) yields \emph{reweight} for covariate-only shift
($\mathrm{AUC}=0.77$), \emph{retrain} for concept drift ($\mathrm{AUC}=0.51$),
and \emph{monitor} for a shift on unused features ($\mathrm{AUC}=0.92$).

\section{Limitations}
\begin{itemize}[leftmargin=2em]
\item MMD is $O(N^2)$; large axes need linear-time MMD or subsampling.
\item Aggregation tends to make detectable shifts separable, so label-free
      covariate/concept attribution is often infeasible at the entity level.
\item Concept drift requires (possibly delayed) labels.
\item Permutation validity assumes exchangeability; correlated entities need
      block permutation. Reported post-selection intervals should use FCR
      correction.
\end{itemize}

\begin{thebibliography}{9}
\bibitem{gretton} A.\ Gretton et al. \emph{A Kernel Two-Sample Test}. JMLR, 2012.
\bibitem{szekely} G.\ Sz\'ekely, M.\ Rizzo. \emph{DISCO analysis / energy distance}. Ann.\ Appl.\ Stat., 2010.
\bibitem{shapley} L.\ Shapley. \emph{A value for $n$-person games}. 1953.
\bibitem{mb} N.\ Meinshausen, P.\ B\"uhlmann. \emph{Stability Selection}. JRSS-B, 2010.
\bibitem{wy} P.\ Westfall, S.\ Young. \emph{Resampling-Based Multiple Testing} (step-down max-T). Wiley, 1993.
\bibitem{bb} Y.\ Benjamini, M.\ Bogomolov. \emph{Selective inference on hierarchies of hypotheses}. JRSS-B, 2014.
\bibitem{chwialkowski} K.\ Chwialkowski et al. \emph{A wild bootstrap for degenerate kernel tests}. NeurIPS, 2014.
\bibitem{niewager} X.\ Nie, S.\ Wager. \emph{Quasi-oracle estimation of heterogeneous treatment effects}. Biometrika, 2021.
\bibitem{garg} S.\ Garg et al. \emph{Leveraging unlabeled data to predict out-of-distribution performance (ATC)}. ICLR, 2022.
\end{thebibliography}

\end{document}
